% !TeX root = ../../../../../main.tex
% chapters/sections/chapter3/subsections/producers/overview.tex

\subsection{生産者の概要}
\label{sec:chapter3_producers_overview}
前節\ref{sec:chapter3_households}で述べたように、本モデルは家計が生産者でもあるヨーマン・ファーマ・モデルである。
家計が直面する複雑な期待生涯効用の最大化問題は「価格以外の変数の最適化問題（家計としての問題）」と「価格の最適化問題（生産者としての問題）」の2つに分割して分析することができ、前節では家計としての最適化について論じた。
本節では、分割されたもう一方の問題である価格の最適化（生産者としての最適化）について論じる。
具体的には、価格改定の機会を得た生産者が将来にわたって価格を据え置く可能性を考慮したうえで、自身の生産する財の価格をどのように決定するかを定式化し、その最適化行動から線形化されたフィリップス曲線を導出する過程について述べる。